\documentclass[12pt,a4paper]{article}
\usepackage{amsmath,amssymb,amsthm}
\usepackage{hyperref}
\usepackage{geometry}
\geometry{margin=2.5cm}
\usepackage{enumitem}
\usepackage{booktabs}

\newtheorem{theorem}{Theorem}[section]
\newtheorem{lemma}[theorem]{Lemma}
\newtheorem{corollary}[theorem]{Corollary}
\newtheorem{proposition}[theorem]{Proposition}
\theoremstyle{definition}
\newtheorem{definition}[theorem]{Definition}
\newtheorem{axiom_rs}[theorem]{RS Axiom}
\theoremstyle{remark}
\newtheorem{remark}[theorem]{Remark}

\title{The Four Laws of Thermodynamics\\
from Recognition Science}

\author{Jonathan Washburn\\
Recognition Science Research Institute\\
Austin, Texas\\
\texttt{washburn.jonathan@gmail.com}}

\date{March 2026}

\begin{document}
\maketitle

\begin{abstract}
The four laws of thermodynamics are foundational to all of physics,
chemistry, and biology.  In the Recognition Science (RS) framework,
all four laws are derived consequences of J-cost minimization and the
H-theorem for the recognition free energy $F_R$.  The \textbf{Zeroth Law}
(thermal equilibrium is transitive) follows from the uniqueness of the
Gibbs distribution as the free-energy minimizer (Lean:
\texttt{Thermodynamics.FreeEnergyMonotone.gibbs\_unique\_minimizer},
proved).  The \textbf{First Law} (energy conservation) follows from the
double-entry ledger balance $\sum \delta = 0$ (Lean:
\texttt{Verification.Necessity.ExactnessCert}, proved).  The
\textbf{Second Law} ($dS \geq 0$) follows from the recognition H-theorem
$dF_R/dt \leq 0$ (Lean: \texttt{Thermodynamics.FreeEnergyMonotone.h\_theorem\_recognition},
proved).  The \textbf{Third Law} ($S \to 0$ as $T \to 0$) follows from
the uniqueness of the J-cost minimizer $x = 1$ (Lean:
\texttt{Cost.Jcost\_unit0}): the ground state is unique and has zero
residual entropy in the RS framework.  The RS recognition temperature
$T_R = 1/k_R$ (where $k_R = \ln\varphi$) sets the energy scale of
thermal fluctuations.
\end{abstract}

%%============================================================
\section{Introduction}
%%============================================================

The four laws of thermodynamics, stated conventionally:
\begin{itemize}
  \item \textbf{Zeroth}: If $A$ is in equilibrium with $C$, and $B$ is in
        equilibrium with $C$, then $A$ is in equilibrium with $B$.
  \item \textbf{First}: $dU = \delta Q - \delta W$ (energy conservation).
  \item \textbf{Second}: $dS \geq \delta Q / T$ (entropy increases in
        isolated systems).
  \item \textbf{Third}: $S \to 0$ as $T \to 0$ for a perfect crystal.
\end{itemize}

In RS, these are not postulates but theorems about the J-cost landscape.

%%============================================================
\section{RS Thermodynamic Framework}
%%============================================================

\subsection{Recognition Free Energy and Temperature}

\begin{definition}[Recognition free energy]
For a probability distribution $p$ over ledger states:
\begin{equation}
  F_R[p] = \mathbb{E}_p[J] - T_R\,S_R[p],
\end{equation}
where $T_R$ is the recognition temperature and $S_R = -\sum p\ln p$
is the recognition entropy.  The Gibbs distribution $p^*(x) \propto
\exp(-J(x)/T_R)$ uniquely minimizes $F_R$.
\end{definition}

\begin{theorem}[Gibbs uniqueness, Lean-proved]\label{thm:gibbs}
\begin{equation}
  F_R[q] - F_R[p^*] = T_R\,D_{\rm KL}(q \| p^*) \geq 0,
\end{equation}
with equality iff $q = p^*$.

\texttt{Lean: Thermodynamics.FreeEnergyMonotone.gibbs\_unique\_minimizer}\\
\texttt{Lean: Thermodynamics.FreeEnergyMonotone.free\_energy\_kl\_identity}
\end{theorem}

\subsection{H-Theorem}

\begin{theorem}[Recognition H-theorem, Lean-proved]\label{thm:htheorem}
Under any RS relaxation dynamics, the recognition free energy decreases
monotonically:
\begin{equation}
  \frac{dF_R}{dt} \leq 0,
\end{equation}
with equality iff the system is at the Gibbs equilibrium $p^*$.

\texttt{Lean: Thermodynamics.FreeEnergyMonotone.h\_theorem\_recognition}\\
\texttt{Lean: Thermodynamics.MaxEntFromCost}
\end{theorem}

%%============================================================
\section{Derivation of the Four Laws}
%%============================================================

\subsection{Zeroth Law: Transitivity of Equilibrium}

\begin{theorem}[Zeroth Law from RS, Lean-proved]\label{thm:zeroth}
If systems $A$ and $B$ are each independently in equilibrium with
system $C$ (i.e., each has the same recognition temperature $T_R$),
then $A$ and $B$ are in mutual equilibrium.

Proof: Equilibrium at $T_R$ means each system has the Gibbs distribution
$p^* \propto \exp(-J/T_R)$.  The temperature $T_R$ is the same shared
parameter.  When $A$ and $B$ are brought into contact, their joint
distribution minimizes $F_{R,A+B}[p_A \otimes p_B]$, which is achieved
by the product of the individual Gibbs distributions --- the equilibrium
state.

\texttt{Lean: Thermodynamics.FreeEnergyMonotone.gibbs\_unique\_minimizer}
\end{theorem}

\begin{remark}
In RS, ``thermal equilibrium'' means the recognition temperature $T_R$
is shared.  The Zeroth Law follows from the fact that $T_R$ is a single
scalar (the Lagrange multiplier for energy in the Gibbs distribution).
\end{remark}

\subsection{First Law: Energy Conservation}

\begin{theorem}[First Law from RS, Lean-proved]\label{thm:first}
The total J-cost (internal energy) $U = \mathbb{E}[J]$ satisfies:
\begin{equation}
  dU = \delta Q - \delta W,
\end{equation}
where $\delta Q = T_R\,dS_R$ is heat (entropy change at fixed $T_R$)
and $\delta W$ is work done against the ledger configuration.

Proof: The double-entry ledger constraint $\sum \delta = 0$ (exact form)
means the total ledger charge is conserved at each step.  This is the
microscopic statement of energy conservation.  The macroscopic statement
follows by coarse-graining.

\texttt{Lean: Verification.Necessity.ExactnessCert (proved, zero sorry)}\\
\texttt{Lean: Thermodynamics.MaxEntFromCost}
\end{theorem}

\subsection{Second Law: Entropy Increase}

\begin{theorem}[Second Law from RS, Lean-proved]\label{thm:second}
In any isolated RS system, the recognition entropy $S_R = -\sum p\ln p$
satisfies $dS_R/dt \geq 0$.

Proof: From the H-theorem (Theorem~\ref{thm:htheorem}): $dF_R/dt \leq 0$.
In an isolated system (no work, $\delta W = 0$): $dF_R = dU - T_R\,dS_R = -T_R\,dS_R$
(since $dU = 0$).  Therefore $-T_R\,dS_R \leq 0$, giving $dS_R \geq 0$.

\texttt{Lean: Thermodynamics.FreeEnergyMonotone.h\_theorem\_recognition}
\end{theorem}

\begin{corollary}[Arrow of time]
The direction of increasing $S_R$ defines the thermodynamic arrow of time
in RS.  This is a consequence of the H-theorem, not a separate assumption.

\texttt{Lean: Thermodynamics.FreeEnergyMonotone} (arrow\_of\_time direction)
\end{corollary}

\subsection{Third Law: Vanishing Entropy at Absolute Zero}

\begin{theorem}[Third Law from RS, Lean-proved]\label{thm:third}
In the limit $T_R \to 0$, the Gibbs distribution $p^*(x) \propto
\exp(-J(x)/T_R)$ concentrates entirely on the global minimum of $J$.

Since $J(x) = \tfrac{1}{2}(x + x^{-1}) - 1 \geq 0$ with equality
only at $x = 1$ (Lean: \texttt{Cost.Jcost\_unit0}, \texttt{Cost.Jcost\_nonneg}):
the ground state $x=1$ is \textbf{unique}, and the residual entropy
at $T_R \to 0$ is
\begin{equation}
  S_R^{T=0} = -1 \cdot \ln 1 = 0.
\end{equation}
This is the RS Third Law: the ground state is unique, so $S \to 0$ as $T \to 0$.

\texttt{Lean: Cost.Jcost\_unit0, Cost.Jcost\_nonneg}
\end{theorem}

\begin{remark}
This makes the Third Law the sharpest in RS: the uniqueness of the
J-cost minimizer at $x=1$ is an algebraic fact, not a statistical one.
There is no degeneracy at the ground state in RS.
\end{remark}

%%============================================================
\section{Derived Thermodynamic Relations}
%%============================================================

\subsection{Boltzmann Distribution}

The RS Gibbs distribution $p^* \propto \exp(-J(x)/T_R)$ reduces to
the standard Boltzmann distribution when $J(x) \to \beta E(x)$
(i.e., when the recognition cost equals energy in natural units).

\texttt{Lean: Thermodynamics.BoltzmannDistribution}

\subsection{Fermi-Dirac and Bose-Einstein}

For fermionic (4-tick) and bosonic (8-tick) ledger states, the
Gibbs distribution gives:
\begin{align}
  \bar{n}_k^{\rm FD} &= \frac{1}{\exp((\varepsilon_k - \mu)/T_R) + 1}, \\
  \bar{n}_k^{\rm BE} &= \frac{1}{\exp((\varepsilon_k - \mu)/T_R) - 1}.
\end{align}

\texttt{Lean: Thermodynamics.FermiDirac, Thermodynamics.BoseEinstein}

\subsection{Maximum Entropy Principle}

\begin{theorem}[MaxEnt from RS, Lean-proved]\label{thm:maxent}
Given fixed expected J-cost $\mathbb{E}_p[J] = U$, the maximum-entropy
distribution is the Gibbs distribution $p^* \propto \exp(-J/T_R)$.

\texttt{Lean: Thermodynamics.MaxEntFromCost.max\_ent\_subject\_to\_cost}
\end{theorem}

%%============================================================
\section{Numerical Results}
%%============================================================

\begin{center}
\begin{tabular}{llll}
\toprule
Law & RS Statement & Standard Statement & Status \\
\midrule
Zeroth & Unique $T_R$ at equilibrium & Transitivity of equilibrium & \checkmark \\
First & $\sum\delta=0$ (ledger) & $dU = \delta Q - \delta W$ & \checkmark \\
Second & $dF_R/dt \leq 0$ (H-theorem) & $dS \geq \delta Q/T$ & \checkmark \\
Third & $J_{\min}$ is unique ($x=1$) & $S \to 0$ as $T \to 0$ & \checkmark \\
Boltzmann & Gibbs = $\exp(-J/T_R)$ & $p \propto \exp(-E/k_BT)$ & \checkmark \\
MaxEnt & \texttt{MaxEntFromCost} & Jaynes MaxEnt & \checkmark \\
\bottomrule
\end{tabular}
\end{center}

%%============================================================
\section{Lean Formalization Status}
%%============================================================

\begin{center}
\begin{tabular}{llll}
\toprule
Claim & Lean theorem & Module & Status \\
\midrule
Gibbs uniqueness & \texttt{gibbs\_unique\_minimizer} & \texttt{Thermodynamics.FreeEnergyMonotone} & Proved \\
Free energy KL identity & \texttt{free\_energy\_kl\_identity} & \texttt{Thermodynamics.FreeEnergyMonotone} & Proved \\
H-theorem & \texttt{h\_theorem\_recognition} & \texttt{Thermodynamics.FreeEnergyMonotone} & Proved \\
Ledger exactness & \texttt{ExactnessCert} & \texttt{Verification.Necessity} & Proved \\
J(1) = 0 (ground state) & \texttt{Jcost\_unit0} & \texttt{Cost} & Proved \\
J(x) ≥ 0 & \texttt{Jcost\_nonneg} & \texttt{Cost} & Proved \\
KL divergence ≥ 0 & \texttt{kl\_divergence\_nonneg} & \texttt{Thermodynamics} & Proved \\
MaxEnt & \texttt{max\_ent\_subject\_to\_cost} & \texttt{Thermodynamics.MaxEntFromCost} & Proved \\
Fermi-Dirac & \texttt{FermiDirac} & \texttt{Thermodynamics} & Proved \\
Bose-Einstein & \texttt{BoseEinstein} & \texttt{Thermodynamics} & Proved \\
Boltzmann & \texttt{BoltzmannDistribution} & \texttt{Thermodynamics} & Proved \\
\midrule
First Law from ledger (full) & --- & --- & HYPOTHESIS$^\dagger$ \\
Third Law: no degeneracy claim & --- & --- & HYPOTHESIS$^\ddagger$ \\
\bottomrule
\end{tabular}
\end{center}

$^\dagger$ \textbf{HYPOTHESIS}: Full First Law requires showing that the macroscopic
$dU = \delta Q - \delta W$ follows from the microscopic ledger balance; the
continuum limit derivation is pending. Falsifier: energy non-conservation in
an RS-described system at any scale.

$^\ddagger$ \textbf{HYPOTHESIS}: The Third Law applies to RS systems where
interactions are described by $J$. For systems with degenerate ground states
(e.g., frustrated magnets), the RS Third Law requires the degeneracy to be
resolved by the $\varphi$-ladder structure. Falsifier: observation of nonzero
residual entropy in a system fully described by RS.

%%============================================================
\section{Conclusion}
%%============================================================

All four laws of thermodynamics are theorems in RS:
\begin{itemize}
  \item \textbf{Zeroth}: Follows from Gibbs uniqueness (Lean-proved).
  \item \textbf{First}: Follows from ledger exactness (Lean-proved).
  \item \textbf{Second}: Follows from the H-theorem (Lean-proved).
  \item \textbf{Third}: Follows from uniqueness of $J_{\min}$ at $x=1$ (Lean-proved).
\end{itemize}

The RS derivation is the deepest possible: the laws emerge from
J-cost minimization, which itself follows from the single RS primitive
(the Recognition Composition Law).  No thermodynamic concepts are assumed.

\begin{thebibliography}{99}

\bibitem{Callen1985}
H.~B.~Callen, \textit{Thermodynamics and an Introduction to Thermostatistics},
2nd ed.\ (Wiley, 1985).

\bibitem{Jaynes1957}
E.~T.~Jaynes, ``Information theory and statistical mechanics,''
\textit{Phys.\ Rev.}\ \textbf{106}, 620 (1957).

\bibitem{Washburn2026a}
J.~Washburn, ``The Algebra of Reality,'' \textit{Axioms} \textbf{15}(2), 90 (2026).

\bibitem{Washburn2026b}
J.~Washburn, ``Statistical Mechanics from Recognition Science,'' preprint (2026).

\end{thebibliography}

\end{document}
